\documentclass[11pt]{article}

% ---------- Packages (all available by default on Overleaf) ----------
\usepackage[a4paper,margin=2.2cm]{geometry}
\usepackage{amsmath,amssymb}
\usepackage{xcolor}
\usepackage{listings}
\usepackage{enumitem}
\usepackage{booktabs}
\usepackage{longtable}
\usepackage{array}
\usepackage{hyperref}
\usepackage{titlesec}
\usepackage{tcolorbox}
\tcbuselibrary{skins,breakable}

\hypersetup{colorlinks=true,linkcolor=blue!50!black,urlcolor=blue!50!black,citecolor=blue!50!black}
\titleformat{\section}{\Large\bfseries\color{blue!35!black}}{\thesection}{0.6em}{}
\titleformat{\subsection}{\large\bfseries\color{blue!25!black}}{\thesubsection}{0.6em}{}
\setlength{\parskip}{2pt}

% ---------- Lean4 (pseudo-code) style ----------
% Code blocks use an ASCII-transliterated form of Lean 4 / Mathlib syntax
% (e.g. "->*" for the unicode arrow of a MonoidHom, "forall"/"exists"/"in"
% for the quantifier and membership glyphs) so the document compiles
% cleanly under plain pdflatex without any special font packages.
\definecolor{codebg}{RGB}{246,248,250}
\definecolor{codekw}{RGB}{175,0,120}
\definecolor{codecomment}{RGB}{110,110,110}
\lstdefinelanguage{Lean4}{
  morekeywords={class,structure,extends,instance,where,def,lemma,example,noncomputable,
    by,intro,rintro,rcases,use,constructor,exact,apply,dsimp,rw,simp,symm,trans,
    fun,forall,exists,in,or,and,calc,inferInstance},
  sensitive=true,
  morecomment=[l]{--},
  morecomment=[s]{/-}{-/},
}
\lstset{
  backgroundcolor=\color{codebg},
  basicstyle=\ttfamily\small,
  keywordstyle=\color{codekw}\bfseries,
  commentstyle=\color{codecomment}\itshape,
  showstringspaces=false,
  frame=single,
  rulecolor=\color{gray!40},
  framesep=4pt,
  breaklines=true,
  language=Lean4,
  xleftmargin=8pt,
  xrightmargin=8pt,
}

% ---------- Colored boxes ----------
\newtcolorbox{defbox}[2]{
  colback=gray!4, colframe=#1!60!black,
  colbacktitle=#1!60!black, coltitle=white,
  fonttitle=\bfseries\ttfamily, title={#2},
  boxrule=0.5pt, arc=2pt, left=6pt, right=6pt, top=3pt, bottom=3pt,
  breakable
}
\newtcolorbox{takeaway}[1][]{
  colback=blue!5!white, colframe=blue!40!black,
  boxrule=0.6pt, arc=2pt, left=6pt, right=6pt, top=4pt, bottom=4pt,
  breakable, title={\textbf{Key Idea}}, fonttitle=\bfseries, #1
}
\newtcolorbox{auditnote}[1][]{
  colback=red!4!white, colframe=red!45!black,
  boxrule=0.6pt, arc=2pt, left=6pt, right=6pt, top=4pt, bottom=4pt,
  breakable, title={\textbf{Audit Note}}, fonttitle=\bfseries, #1
}

\title{\vspace{-1.5cm}\textbf{Chapter 9 Summary}\\[2pt]\Large Groups, Rings, and Their Quotients in Lean 4}
\author{Reading notes}
\date{}

\begin{document}
\maketitle
\vspace{-0.6cm}

\begin{abstract}
\noindent This chapter moves from the abstract algebraic classes of Chapter~8 to concrete Mathlib objects: group and ring homomorphisms, subgroups and ideals, kernels and images, quotient groups/rings, and the isomorphism theorems that relate them, closing with a short tour of polynomial rings. This note restates the chapter's results with corrected mathematical statements and completed proof terms where the source notes contained typos or ambiguous phrasing; every such correction is flagged in an \textcolor{red!45!black}{\textbf{Audit Note}} box.
\end{abstract}
\vspace{-0.3cm}
\tableofcontents

% =====================================================================
\section*{Notation Used in This Note}
\addcontentsline{toc}{section}{Notation Used in This Note}
% =====================================================================
\renewcommand{\arraystretch}{1.25}
\begin{center}
\begin{tabular}{@{}p{3.4cm}p{5.6cm}p{5.6cm}@{}}
\toprule
\textbf{Written here} & \textbf{Meaning} & \textbf{Lean / Mathlib notation} \\
\midrule
$f : G \to H$ & group homomorphism & \verb|f : G ->* H| \\
$G \cong H$ & group isomorphism & \verb|G <=>* H| \\
$f : R \to S$ & ring homomorphism & \verb|f : R ->+* S| \\
$R \cong S$ & ring isomorphism & \verb|R <=>+* S| \\
$G/H$, $R/I$ & quotient group / ring & \verb|G / H|, \verb|R / I| \\
$\ker f$, $\operatorname{im} f$ & kernel, image of a hom. & \verb|f.ker|, \verb|f.range| \\
$R^\times$ & the group of units of $R$ & \verb|Units R| \\
\bottomrule
\end{tabular}
\end{center}

% =====================================================================
\section{Group Homomorphisms, Isomorphisms, and Subgroups}
% =====================================================================

\subsection{Homomorphisms and Isomorphisms}

A \emph{group homomorphism} is a \texttt{Monoid} homomorphism between two groups: a function $f:G\to H$ with $f(xy)=f(x)f(y)$ (identity-preservation then follows automatically).

\begin{defbox}{blue}{Monoid / Group homomorphisms and isomorphisms}
\begin{lstlisting}
-- applying a bundled MonoidHom
example {M N : Type*} [Monoid M] [Monoid N] (x y : M) (f : M ->* N) :
    f (x * y) = f x * f y := f.map_mul x y

-- building a group hom from a raw function + a preservation proof
example {G H : Type*} [Group G] [Group H] (f : G -> H)
    (h : forall x y, f (x * y) = f x * f y) : G ->* H :=
  MonoidHom.mk' f h

-- a bijective group hom is a group isomorphism
noncomputable example {G H : Type*} [Group G] [Group H]
    (f : G ->* H) (h : Function.Bijective f) : G <=>* H :=
  MulEquiv.ofBijective f h
\end{lstlisting}
\end{defbox}

Also recall that \texttt{AddCommGroup} lets us discharge routine additive-group identities automatically, e.g.\ $z+x+(y-z-x)=y$ via the \verb|abel| tactic, the additive-group analogue of \verb|ring|.

\subsection{Subgroups}

A subset $S\subseteq G$ is a subgroup exactly when it is closed under the group's structure:
\begin{enumerate}[leftmargin=1.6em,itemsep=1pt]
  \item it contains the identity ($0$ for \texttt{AddGroup}, $1$ for \texttt{Group});
  \item it is closed under the group operation;
  \item it is closed under inverses ($-a$ resp.\ $a^{-1}$);
  \item its elements coerce back into the ambient group (the carrier is genuinely a subset).
\end{enumerate}

\begin{defbox}{blue}{Example: \texorpdfstring{$\mathbb{Z}$}{Z} as a subgroup of \texorpdfstring{$\mathbb{Q}$}{Q}}
\begin{lstlisting}
example : AddSubgroup Rat where
  -- carrier = { y in Q | exists x in Z, (x : Q) = y }
  carrier := Set.range ((coe) : Int -> Rat)
  add_mem'  := by rintro _ _ (n, rfl) (m, rfl); use n + m; simp
  zero_mem' := by use 0; simp
  neg_mem'  := by rintro _ (n, rfl); use -n; simp
\end{lstlisting}
\end{defbox}

\begin{defbox}{blue}{Example: conjugation preserves the subgroup property}
For $x\in G$ and a subgroup $H\le G$, the \emph{conjugate} $xHx^{-1}=\{xhx^{-1}\mid h\in H\}$ is again a subgroup:
\begin{lstlisting}
def conjugate {G : Type*} [Group G] (x : G) (H : Subgroup G) : Subgroup G where
  carrier := {a : G | exists h, h in H and a = x * h * x^{-1}}
  one_mem' := by
    dsimp; use 1
    constructor
    . exact H.one_mem
    . symm; rw [mul_one, mul_inv_cancel]
  inv_mem' := by
    dsimp; rintro x1 (h, h_in, rfl)
    use h^{-1}, H.inv_mem h_in; group
  mul_mem' := by
    dsimp; rintro a b (ha, ha_in, rfl) (hb, hb_in, rfl)
    use ha * hb, mul_mem ha_in hb_in; group
\end{lstlisting}
\end{defbox}

% =====================================================================
\section{Kernels, Images, and the Subgroup Lattice}
% =====================================================================

For a homomorphism $f:G\to H$, its \emph{kernel} $\ker f=\{g\in G\mid f(g)=1\}$ and \emph{image} $\operatorname{im}f=\{h\in H\mid \exists g,\,f(g)=h\}$ are subgroups of $G$ and $H$ respectively:

\begin{lstlisting}
example {G H : Type*} [Group G] [Group H] (f : G ->* H) (g : G) :
    g in MonoidHom.ker f <-> f g = 1 := f.mem_ker

example {G H : Type*} [Group G] [Group H] (f : G ->* H) (h : H) :
    h in MonoidHom.range f <-> exists g : G, f g = h := f.mem_range
\end{lstlisting}

More generally, a homomorphism $f:G\to H$ pushes subgroups of $G$ forward (\texttt{map}) and pulls subgroups of $H$ back (\texttt{comap}); both operations are monotone and respect composition, so together they form a Galois-connection-like pair between the subgroup lattices of $G$ and $H$:

\renewcommand{\arraystretch}{1.3}
\begin{center}
\begin{tabular}{@{}p{2.6cm}p{6.4cm}p{5cm}@{}}
\toprule
\textbf{Operation} & \textbf{Statement} & \textbf{Lean lemma / term} \\
\midrule
Map & $G'\le G,\ f:G\to H \ \Rightarrow\ f(G')\le H$ & \verb|Subgroup.map f G'| \\
Comap & $H'\le H,\ f:G\to H \ \Rightarrow\ f^{-1}(H')\le G$ & \verb|Subgroup.comap f H'| \\
Map monotone & $S\le T \ \Rightarrow\ \mathrm{map}\,f\,S \le \mathrm{map}\,f\,T$ & unfolds via \verb|mem_map| \\
Comap monotone & $S\le T \ \Rightarrow\ \mathrm{comap}\,f\,S \le \mathrm{comap}\,f\,T$ & unfolds via \verb|mem_comap| \\
Map composition & $(\psi\circ\varphi)(S) = \psi(\varphi(S))$ & \verb|mem_map| + \verb|ext| \\
Comap composition & $(\psi\circ\varphi)^{-1}(U) = \varphi^{-1}(\psi^{-1}(U))$ & \verb|mem_comap| + \verb|ext| \\
\bottomrule
\end{tabular}
\end{center}

% =====================================================================
\section{Structure Theorems: Lagrange and Sylow}
% =====================================================================

\begin{defbox}{violet}{Lagrange's Theorem}
\begin{lstlisting}
example {G : Type*} [Group G] (G' : Subgroup G) :
    Nat.card G' | Nat.card G := by
  refine (G'.index, ?_)
  rw [mul_comm G'.index]
  exact G'.index_mul_card.symm
\end{lstlisting}
$|G'|$ divides $|G|$; the witness is the \emph{index} $[G:G']$, since \verb|index_mul_card| gives $[G:G']\cdot|G'|=|G|$.
\end{defbox}

\begin{defbox}{violet}{Sylow's Theorem}
\begin{lstlisting}
example {G : Type*} [Group G] [Finite G] (p : Nat) {n : Nat} [Fact p.Prime]
    (hdvd : p ^ n | Nat.card G) : exists K : Subgroup G, Nat.card K = p ^ n :=
  Sylow.exists_subgroup_card_pow_prime p hdvd
\end{lstlisting}
If $p^n$ divides $|G|$ for a prime $p$, a subgroup of order exactly $p^n$ exists.
\end{defbox}

\begin{auditnote}
The source writes the Lagrange witness as a single term using Lean's infix rewrite-substitution combinator, informally $h \triangleright e$ (meaning: rewrite $e$ using the equation $h$, i.e.\ the term-mode equivalent of \verb|rw [h] at e|). The listing above spells this out as an equivalent short tactic proof (\verb|rw|, then \verb|exact|) to keep the document self-contained.
\end{auditnote}

% =====================================================================
\section{Quotient Groups and Isomorphism Theorems}
% =====================================================================

\subsection{Quotient by a Normal Subgroup}

$H\trianglelefteq G$ (every conjugate of $H$ stays inside $H$) is exactly the condition needed for the coset space $G/H$ to inherit a group structure from $G$:

\begin{lstlisting}
example {G : Type*} [Group G] (H : Subgroup G) [H.Normal] : Group (G / H) :=
  inferInstance
\end{lstlisting}

\emph{Proof sketch.} Normality of $H$ means $\forall x\in G,\ h\in H,\ xhx^{-1}\in H$. Multiplication on $G/H$ is defined by $(a_1H)(a_2H) := (a_1a_2)H$; well-definedness needs
\[
a_1H\cdot a_2H \;=\; a_1a_2\,(a_2^{-1}Ha_2)\,H \;=\; a_1a_2H,
\]
which holds precisely because normality gives $a_2^{-1}Ha_2\subseteq H$.

\subsection{The First Isomorphism Theorem}

For $\varphi:G\to M$, the induced map $G/\ker\varphi \to \operatorname{im}\varphi$ is an isomorphism:

\begin{lstlisting}
example {G M : Type*} [Group G] [Group M] (phi : G ->* M) :
    G / MonoidHom.ker phi <=>* MonoidHom.range phi :=
  QuotientGroup.quotientKerEquivRange phi
\end{lstlisting}

\begin{auditnote}
The source states the target type of this example as \verb|G / ker phi ->* range phi| (a mere homomorphism, \verb|->*|). Since \verb|QuotientGroup.quotientKerEquivRange| constructs a bijective correspondence and the surrounding text calls this fact ``Group Equivalence,'' the type should be the isomorphism arrow \verb|<=>*|, as corrected above.
\end{auditnote}

\subsection{Order of a Quotient Group}

If $|G|=|H|\cdot|K|$ then $|G/H|=|K|$, by cancelling $|H|$ out of $|G/H|\cdot|H| = |G| = |K|\cdot|H|$:

\begin{lstlisting}
lemma aux_card_eq [Finite G] (h' : Nat.card G = Nat.card H * Nat.card K) :
    Nat.card (G / H) = Nat.card K := by
  have : Nat.card (G / H) * Nat.card H = Nat.card K * Nat.card H := calc
    Nat.card (G / H) * Nat.card H = Nat.card G       := by rw [<- index_eq_card H, index_mul_card H]
    _                              = Nat.card K * Nat.card H := by rw [h', mul_comm]
  exact Nat.eq_of_mul_eq_mul_right Nat.card_pos this
\end{lstlisting}
using \verb|index_eq_card H| ($[G:H]=|G/H|$) and \verb|index_mul_card H| ($[G:H]\cdot|H|=|G|$).

\subsection{A Direct-Product Recognition Theorem}

The source's final three results are labeled ``First/Second/Third Group Isomorphism Theorem,'' but they do not correspond to the classically-numbered isomorphism theorems (whose usual statements are $G/\ker\varphi\cong\operatorname{im}\varphi$; the diamond theorem $HN/N\cong H/(H\cap N)$; and $(G/N)/(M/N)\cong G/M$). They are instead the three steps of a single argument --- under the hypotheses $H,K\le G$, $|G|=|H|\cdot|K|$, and $H\cap K=\{1\}$ --- that $G$ is (isomorphic to) the internal direct product $H\times K$.

\begin{auditnote}
The source also describes the hypothesis as ``$H,G$ are disjoint,'' which cannot be literal (a subgroup always contains $G$'s identity, so $H$ and $G$ trivially intersect at $1$, and $H\subseteq G$ regardless). The intended hypothesis, consistent with the Lean proof terms, is that \emph{$H$ and $K$ intersect trivially}, $H\cap K=\{1\}$; the phrasing below uses that reading.
\end{auditnote}

\begin{defbox}{teal}{Step 1: \texorpdfstring{$K\cong G/H$}{K iso G/H}}
\begin{lstlisting}
def iso1 : K <=>* G / H := by
  apply MulEquiv.ofBijective ((QuotientGroup.mk' H).restrict K)
  rw [Nat.bijective_iff_injective_and_card]
  constructor
  . rw [<- ker_eq_bot_iff, (QuotientGroup.mk' H).ker_restrict K]; simp [h]
  . symm; exact aux_card_eq h'
\end{lstlisting}
\verb|QuotientGroup.mk' H| is the canonical hom $G\to G/H$; restricting it to $K$ is injective exactly when $K\cap H=\{1\}$, and \verb|Nat.bijective_iff_injective_and_card| upgrades injective + matching cardinalities to bijective.
\end{defbox}

\begin{defbox}{teal}{Step 2: \texorpdfstring{$G\cong (G/K)\times (G/H)$}{G iso (G/K) x (G/H)}}
\begin{lstlisting}
def iso2 : G <=>* (G / K) * (G / H) := by
  -- the pair of canonical projections (G -> G/K, G -> G/H) combined
  apply MulEquiv.ofBijective <| (QuotientGroup.mk' K).prod (QuotientGroup.mk' H)
  rw [Nat.bijective_iff_injective_and_card]
  constructor
  . rw [<- ker_eq_bot_iff, ker_prod]; simp [h.symm.eq_bot]
  . rw [Nat.card_prod, aux_card_eq h', aux_card_eq (mul_comm (Nat.card H) _ substEq h'), h']
\end{lstlisting}
The product of the two canonical projections is injective because the kernels' intersection is trivial, and the cardinalities match by hypothesis. (\verb|substEq| stands in for Lean's built-in term-level rewrite-substitution combinator, spelled out as an identifier here to keep the listing plain ASCII.)
\end{defbox}

\begin{defbox}{teal}{Step 3: the recognition theorem \texorpdfstring{$G\cong H\times K$}{G iso H x K}}
\begin{lstlisting}
def finalIso : G <=>* H * K :=
  (iso2 h h').trans (((iso1 h.symm (mul_comm (Nat.card H) _ substEq h')).prodCongr
    (iso1 h h')).symm)
\end{lstlisting}
Chain: Step~2 reduces $G\cong H\times K$ to $G\cong (G/K)\times(G/H)$; Step~1 (applied twice, with roles of $H,K$ swapped) turns $(G/K)\times(G/H)$ into $H\times K$; composing gives $G\cong H\times K$.
\end{defbox}

% =====================================================================
\section{Ring Homomorphisms and Units}
% =====================================================================

\begin{defbox}{orange}{Facts 1--4: units of a ring}
\begin{lstlisting}
-- Fact 1: the only units of Z are 1 and -1
example (x : Units Int) : x = 1 or x = -1 := Int.units_eq_one_or x

-- Fact 2: a unit times its inverse is 1 (Monoid is multiplicative by default)
example {M : Type*} [Monoid M] (x : Units M) : (x : M) * x^{-1} = 1 := Units.mul_inv x

-- Fact 3: the units of a monoid form a group
example {M : Type*} [Monoid M] : Group (Units M) := inferInstance

-- Fact 4: a ring hom induces a hom between unit groups
example {R S : Type*} [Ring R] [Ring S] (f : R ->+* S) : Units R ->* Units S :=
  Units.map f
\end{lstlisting}
\end{defbox}

% =====================================================================
\section{Ideals and Quotient Rings}
% =====================================================================

\begin{defbox}{orange}{Facts 5--8: the quotient ring \texorpdfstring{$R/I$}{R/I}}
\begin{lstlisting}
-- Fact 5: every ideal gives a canonical quotient ring hom
example {R : Type*} [CommRing R] (I : Ideal R) : R ->+* R / I :=
  Ideal.Quotient.mk I

-- Fact 6: elements of I are exactly the elements sent to 0
example {R : Type*} [CommRing R] {a : R} {I : Ideal R} :
    Ideal.Quotient.mk I a = 0 <-> a in I :=
  Ideal.Quotient.eq_zero_iff_mem

-- Fact 7: universal property -- f factors through R / I whenever I <= ker f
example {R S : Type*} [CommRing R] [CommRing S] (I : Ideal R) (f : R ->+* S)
    (H : I <= RingHom.ker f) : R / I ->+* S :=
  Ideal.Quotient.lift I f H

-- Fact 8: the (ring) First Isomorphism Theorem
example {R S : Type*} [CommRing R] [CommRing S] (f : R ->+* S) :
    R / RingHom.ker f <=>+* f.range :=
  RingHom.quotientKerEquivRange f
\end{lstlisting}
Fact~5 is the ring analogue of \verb|QuotientGroup.mk'|; Fact~7 is the universal property that makes $R/I$ the ``most efficient'' quotient satisfying $I\subseteq\ker f$; Fact~8 is the ring version of the group isomorphism theorem from Section~5.2.
\end{defbox}

\begin{defbox}{orange}{Facts 9--14: ideal arithmetic and quotient maps}
For ideals $I,J$ of a commutative ring $R$:
\end{defbox}
\renewcommand{\arraystretch}{1.3}
\begin{center}
\begin{tabular}{@{}p{1.3cm}p{7cm}p{5.5cm}@{}}
\toprule
\textbf{Fact} & \textbf{Statement} & \textbf{Lean lemma / term} \\
\midrule
9 & $x\in I+J \iff \exists a\in I,\, b\in J,\ a+b=x$ & \verb|Submodule.mem_sup| \\
10 & $I\cdot J \subseteq J$ & \verb|Ideal.mul_le_left| \\
11 & $I\cdot J \subseteq I$ & \verb|Ideal.mul_le_right| \\
12 & $I\cdot J \subseteq I\sqcap J$ & \verb|Ideal.mul_le_inf| \\
13 & $I\subseteq f^{-1}(J) \Rightarrow$ exists a lift $R/I\to S/J$ & \verb|Ideal.quotientMap| \\
14 & $I=J \Rightarrow R/I \cong R/J$ & \verb|Ideal.quotEquivOfEq| \\
\bottomrule
\end{tabular}
\end{center}

\begin{takeaway}
Facts 10--12 all follow the same pattern: since $\forall x\in R,\ x\cdot J\subseteq J$ (resp.\ $I\cdot x\subseteq I$) by definition of ideal multiplication, restricting $x$ to range over the smaller set $I$ (resp.\ $J$) only shrinks the union, giving both containments and hence containment in the intersection $I\sqcap J$.
\end{takeaway}

% =====================================================================
\section{Polynomial Rings}
% =====================================================================

\begin{defbox}{orange}{Fact 15: arithmetic, the embedding \texorpdfstring{$C$}{C}, and coefficients in \texorpdfstring{$R[X]$}{R[X]}}
\begin{lstlisting}
-- C r is the constant polynomial with value r; C : R ->+* R[X] is a ring hom
example {R : Type*} [CommRing R] (r : R) : (X + C r) * (X - C r) = X ^ 2 - C (r ^ 2) := by
  rw [C.map_pow]; ring

example {R : Type*} [CommRing R] (r : R) : (C r).coeff 0 = r := by simp

example {R : Type*} [CommRing R] :
    (X ^ 2 + 2 * X + C 3 : Poly R).coeff 1 = 2 := by simp
\end{lstlisting}
Because $C$ is a ring homomorphism, $(C r)^2 = C(r^2)$, which is what finishes the first identity once \verb|ring| handles the remaining polynomial algebra. The coefficient lemmas show \verb|coeff| reading off $r$ from the constant polynomial $C r$, and the coefficient of $X^1$ from $X^2+2X+C\,3$.
\end{defbox}

\begin{auditnote}
Two issues in the source were corrected above: (1) the header claimed these examples show that ``\verb|R[X]| is also a commutative ring,'' but \verb|CommRing (R[X])| is an existing Mathlib instance, not something these snippets establish --- they instead illustrate arithmetic, the embedding $C$, and \verb|coeff|, so the description here was reworded to match; (2) the proof of the first identity ended with the incomplete tactic \verb|rin|, completed above as \verb|ring| (the tactic needed to close the remaining polynomial-ring equality).
\end{auditnote}

\renewcommand{\arraystretch}{1.3}
\begin{center}
\begin{tabular}{@{}p{1.3cm}p{7cm}p{5.5cm}@{}}
\toprule
\textbf{Fact} & \textbf{Statement} & \textbf{Lean lemma / term} \\
\midrule
16 & $\deg(pq) = \deg p + \deg q$ over a semiring with no zero divisors & \verb|Polynomial.degree_mul| \\
17 & $\deg(p\circ q) = \deg p \cdot \deg q$ (natural-number degree) & \verb|Polynomial.natDegree_comp| \\
18 & evaluating $p$ at $x\in R$ gives an element of $R$ & \verb|P.eval x| \\
19 & $r$ is a root of $p$ $\iff$ $p(r)=0$ & \verb|Iff.rfl| \\
\bottomrule
\end{tabular}
\end{center}

\begin{auditnote}
Fact 19's source proof term was written \verb|Iff.rf|, an incomplete/typo'd form of \verb|Iff.rfl| (the proof that \verb|IsRoot P r| unfolds definitionally to \verb|P.eval r = 0|); corrected above.
\end{auditnote}

\begin{defbox}{orange}{Facts 20--21: roots of \texorpdfstring{$(X-C r)^n$}{(X-Cr)\^{}n} over an integral domain}
\begin{lstlisting}
-- Fact 20: (X - C r) has the single root r
example {R : Type*} [CommRing R] [IsDomain R] (r : R) :
    (X - C r).roots = {r} := roots_X_sub_C r

-- Fact 21: (X - C r)^n has root r with multiplicity n
example {R : Type*} [CommRing R] [IsDomain R] (r : R) (n : Nat) :
    ((X - C r) ^ n).roots = nsmul n ({r} : Multiset R) := by simp
\end{lstlisting}
\verb|IsDomain R| (commutative ring, no zero divisors) is required because \verb|.roots| relies on the fact that a nonzero degree-$n$ polynomial has at most $n$ roots \emph{in an integral domain} --- this bound fails over rings with zero divisors. In Fact~21, \verb|nsmul n {r}| is the multiset containing $r$ with multiplicity $n$ (Lean writes this with an infix \texttt{SMul} operator; \verb|nsmul| is used here purely to keep the listing ASCII-only) --- unrelated to the ring's own scalar action.
\end{defbox}

\end{document}
